% =============================================================================
\section{WRP 架构}
\label{sec:wrp}
% =============================================================================

我们的论文体系支持如下论点：

\begin{thesis}[WRP 耦合性]
\label{thesis:wrp}
LLM 推理优化中的三个维度，即 \textbf{Workload}、\textbf{Router} 与 \textbf{Pool}，彼此是耦合的，而不是正交的。若只孤立地优化某一个维度，就会留下大量效率红利没有被利用。最大的收益来自跨维度的联合优化。
\end{thesis}

来自我们自身工作的证据如下：FleetOpt~\cite{fleetopt2026}表明，将路由压缩参数 $\gamma$ 与资源池规模 $(n_s, n_l)$ 共同设计，比在既有集群之上事后加压缩，成本还能再低 3.1--6.4\%（即 \textbf{Router $\times$ Pool} 耦合）。1/W 定律~\cite{onewlaw2026}表明，同一套 GPU 集群会因为所服务的上下文窗口不同而产生 40$\times$ 的能效差异（即 \textbf{Workload $\times$ Pool} 耦合）。AVR~\cite{avr2026}则说明，VLM 工作负载所需的路由信号与纯文本聊天完全不同（即 \textbf{Workload $\times$ Router} 耦合）。

\begin{definition}[WRP 分解]
一个生产级 LLM 推理部署可以由三元组 $(\mathbf{W}, \mathbf{R}, \mathbf{P})$ 来刻画，其中：
\begin{itemize}[nosep]
  \item $\mathbf{W}$：表示\emph{工作负载}，即对请求类型、提示长度、生成长度、会话结构以及模态的分布；
  \item $\mathbf{R}$：表示\emph{路由器}，它将每个请求映射为一个 (model, pool, configuration) 三元组；以及
  \item $\mathbf{P}$：表示\emph{资源池架构}，规定 GPU 类型、资源池边界、解耦拓扑以及 KV-cache 管理方式。
\end{itemize}
优化目标是在 SLO 约束下，对\emph{成本}（每请求 GPU 小时）、\emph{准确率}（任务完成质量）、\emph{延迟}（P99 TTFT、TPOT）和\emph{能耗}（每瓦令牌数）进行加权组合。
\end{definition}
